% GENERATED FILE. Edit canonical lessons, then run npm run generate:books.
% canonical-source-hash: fnv1a64:66db90c8e0134174
% canonical-lessons: GE-C462-unwetter, GE-C462-aufziehen, GE-C462-hagel, GE-C462-hageln, GE-C462-donnern

\chapter{Ein Unwetter}
\label{ch:ge-ein-unwetter}
\hlchaptermodality{462}

\begin{chapteropening}
\textbf{By the end of this chapter:} \emph{I can describe a storm coming: severe weather, a storm building up, hail, hailing and thunder.}

Complete the last lesson of chapter 462: I can describe a storm coming: severe weather, a storm building up, hail, hailing and thunder.
\end{chapteropening}

\section[das Unwetter]{das Unwetter — the severe weather, the violent storm}

\label{lesson:GE-C462-unwetter}



\begin{quote}
Before the new one: say the German for the sunny spell, then the German for to clear up.
\end{quote}

\subsection*{You'll want to know: das Unwetter}
\textbf{das Unwetter} — \textquotedblleft{}the severe weather, the violent storm\textquotedblright{}.

\begin{quote}
Nach dem Unwetter liegen Äste auf der Straße.
\end{quote}

\begin{quote}
\emph{After the storm there are branches lying in the road.}
\end{quote}

Same form in the plural.

\begin{cousinweb}
\emph{Unwetter} is inherited from Middle High German \emph{unweter} 'bad weather' (EtymDB), built from \emph{un-} 'bad, not' and \emph{Wetter} 'weather'.
\end{cousinweb}

\begin{grammarlens}[title={What you've built}]
One more word for telling, explaining and describing at B1.
\end{grammarlens}

\subsection*{Your turn}
\begin{itemize}
\raggedright
  \item \emph{Say it:} \emph{das Unwetter}
  \item \emph{Say it:} \emph{das Unwetter}, once more
  \item \emph{Say it:} a sentence of your own with \emph{das Unwetter}
  \item \emph{Recall:} say the German for the sunny spell, then the German for to clear up, then say \emph{das Unwetter} again
\end{itemize}

\begin{tcolorbox}[breakable,colback=teal!4,colframe=teal!35!black,title={Before you move on}]
What does das Unwetter mean? (\textquotedblleft{}The severe weather, the violent storm\textquotedblright{}.) Say it once more.
\end{tcolorbox}
\section[aufziehen]{aufziehen — to come up, to gather (of a storm); to bring up (a child)}

\label{lesson:GE-C462-aufziehen}



\begin{quote}
Before the new one: say the German for to clear up, then the German for the severe weather.
\end{quote}

\subsection*{You'll want to know: aufziehen}
\textbf{aufziehen} — \textquotedblleft{}to come up, to gather (of a storm); to bring up (a child)\textquotedblright{}.

\begin{quote}
Am Abend zieht ein Gewitter auf.
\end{quote}

\begin{quote}
\emph{In the evening a thunderstorm comes up.}
\end{quote}

Separable; Perfekt with \emph{sein} for weather. \emph{Kinder aufziehen} — to bring up children (with \emph{haben}).

\begin{cousinweb}
\emph{aufziehen} is built from \emph{auf} 'up' and \emph{ziehen} 'to pull, to move', inherited from Proto-Germanic \emph{teuhaną} (EtymDB).
\end{cousinweb}

\begin{grammarlens}[title={What you've built}]
One more word for telling, explaining and describing at B1.
\end{grammarlens}

\subsection*{Your turn}
\begin{itemize}
\raggedright
  \item \emph{Say it:} \emph{aufziehen}
  \item \emph{Say it:} \emph{aufziehen}, once more
  \item \emph{Say it:} a sentence of your own with \emph{aufziehen}
  \item \emph{Recall:} say the German for to clear up, then the German for the severe weather, then say \emph{aufziehen} again
\end{itemize}

\begin{tcolorbox}[breakable,colback=teal!4,colframe=teal!35!black,title={Before you move on}]
What does aufziehen mean? (\textquotedblleft{}To come up, to gather (of a storm); to bring up (a child)\textquotedblright{}.) Say it once more.
\end{tcolorbox}
\section[der Hagel]{der Hagel — the hail}

\label{lesson:GE-C462-hagel}



\begin{quote}
Before the new one: say the German for the severe weather, then the German for to come up.
\end{quote}

\subsection*{You'll want to know: der Hagel}
\textbf{der Hagel} — \textquotedblleft{}the hail\textquotedblright{}.

\begin{quote}
Der Hagel hat viele Autos beschädigt.
\end{quote}

\begin{quote}
\emph{The hail damaged lots of cars.}
\end{quote}

No plural. \emph{Die Hagelkörner} — hailstones.

\begin{cousinweb}
\emph{Hagel} is inherited from Old High German \emph{hagal}, and English \emph{hail} is a cousin. EtymDB lists the verb \emph{hageln} from Middle High German \emph{hagelen}.
\end{cousinweb}

\begin{grammarlens}[title={What you've built}]
One more word for telling, explaining and describing at B1.
\end{grammarlens}

\subsection*{Your turn}
\begin{itemize}
\raggedright
  \item \emph{Say it:} \emph{der Hagel}
  \item \emph{Say it:} \emph{der Hagel}, once more
  \item \emph{Say it:} a sentence of your own with \emph{der Hagel}
  \item \emph{Recall:} say the German for the severe weather, then the German for to come up, then say \emph{der Hagel} again
\end{itemize}

\begin{tcolorbox}[breakable,colback=teal!4,colframe=teal!35!black,title={Before you move on}]
What does der Hagel mean? (\textquotedblleft{}The hail\textquotedblright{}.) Say it once more.
\end{tcolorbox}
\section[hageln]{hageln — to hail}

\label{lesson:GE-C462-hageln}



\begin{quote}
Before the new one: say the German for to come up, then the German for the hail.
\end{quote}

\subsection*{You'll want to know: hageln}
\textbf{hageln} — \textquotedblleft{}to hail\textquotedblright{}.

\begin{quote}
Im Mai hat es plötzlich gehagelt.
\end{quote}

\begin{quote}
\emph{In May it suddenly hailed.}
\end{quote}

With \emph{es}: \emph{es hagelt}, \emph{es hat gehagelt}.

\begin{cousinweb}
\emph{hageln} is inherited from Middle High German \emph{hagelen} (EtymDB), built on \emph{Hagel} 'hail'.
\end{cousinweb}

\begin{grammarlens}[title={What you've built}]
One more word for telling, explaining and describing at B1.
\end{grammarlens}

\subsection*{Your turn}
\begin{itemize}
\raggedright
  \item \emph{Say it:} \emph{hageln}
  \item \emph{Say it:} \emph{hageln}, once more
  \item \emph{Say it:} a sentence of your own with \emph{hageln}
  \item \emph{Recall:} say the German for to come up, then the German for the hail, then say \emph{hageln} again
\end{itemize}

\begin{tcolorbox}[breakable,colback=teal!4,colframe=teal!35!black,title={Before you move on}]
What does hageln mean? (\textquotedblleft{}To hail\textquotedblright{}.) Say it once more.
\end{tcolorbox}
\section[donnern]{donnern — to thunder}

\label{lesson:GE-C462-donnern}



\begin{quote}
Before the new one: say the German for the hail, then the German for to hail.
\end{quote}

\subsection*{You'll want to know: donnern}
\textbf{donnern} — \textquotedblleft{}to thunder\textquotedblright{}.

\begin{quote}
Es blitzt und donnert die ganze Nacht.
\end{quote}

\begin{quote}
\emph{There is thunder and lightning all night.}
\end{quote}

With \emph{es}: \emph{es donnert}. \emph{Der Donner} — the thunder.

\begin{cousinweb}
\emph{donnern} is built on \emph{Donner} 'thunder' (EtymDB lists the verb), inherited from Proto-Germanic \emph{þunraz}; English \emph{thunder} is a cousin.
\end{cousinweb}

\begin{grammarlens}[title={What you've built}]
One more word for telling, explaining and describing at B1.
\end{grammarlens}

\subsection*{Your turn}
\begin{itemize}
\raggedright
  \item \emph{Say it:} \emph{donnern}
  \item \emph{Say it:} \emph{donnern}, once more
  \item \emph{Say it:} a sentence of your own with \emph{donnern}
  \item \emph{Recall:} say the German for the hail, then the German for to hail, then say \emph{donnern} again
\end{itemize}

\begin{tcolorbox}[breakable,colback=teal!4,colframe=teal!35!black,title={Before you move on}]
What does donnern mean? (\textquotedblleft{}To thunder\textquotedblright{}.) Say it once more.
\end{tcolorbox}
